\chapter{Lernaufgabe: Netzwerke}
\textbf{Ziel:} Am Ende soll ein funktionsfähiges Programm stehen, mit dem die Lernenden in Echtzeit Nachrichten austauschen können.

\textbf{Problembezug:} \enquote{Wie kann ich mit einer anderen Person eine Nachricht über das Internet austauschen?}

\begin{mainBox}[title=Komplexe Lernaufgabe: Chat-Anwendung]
    Die \textbf{Haupt-Lernaufgabe} lautet: Entwickeln Sie in Kleingruppen eine \textbf{funktionierende Chat-Anwendung} in Python, die über einen Relay-Server läuft. Am Ende sollen die Gruppen miteinander in Echtzeit Nachrichten austauschen können. Diese Haupt-Lernaufgabe ist in mehrere \textbf{Learning Tasks} unterteilt, die aufeinander aufbauen und jeweils durch unterstützende Informationen, Just-in-time Hilfen und Part-task Practice (Teilübungen) begleitet werden.
\end{mainBox}

\begin{scaffoldBox}[title=Scaffolding der Lernaufgaben]
    Die Unterrichtssequenz besteht aus drei aufeinander aufbauenden Learning Tasks, welche jeweils durch vorausgehende, unterstützende Informationen, Just-in-time Hilfen und Part-task Practice (Teilübungen) begleitet werden. Dieses \textit{Scaffolding} jedes Learning Tasks wird jeweils in einer blauen Box hervorgehoben, ebenso wie das fortschreitende Zurückweichen (\textit{Fading}) der Unterstützung.
\end{scaffoldBox}

\begin{table}[H]
    \centering
    \small
    \renewcommand{\arraystretch}{1.2}
    \begin{tabular}{p{0.33\linewidth}|p{0.33\linewidth}|p{0.33\linewidth}}
        \textbf{Unterstützende Informationen}                                                  & \textbf{Just-in-time Informationen} & \textbf{Part-task Practice (Teilübungen)} \\
        \midrule

        \begin{itemize}
            \item Visualisierung der Client-Server-Architektur (Skript, Slides)
            \item Einführung in grundlegende Begriffe: IP-Adresse, Port, Socket
            \item Analogie: \enquote{Postsystem} (Adresse $\rightarrow$ Empfänger)
            \item Infoblätter / Moodle-Materialien mit Hintergrundwissen
        \end{itemize}                 &
        \begin{itemize}
            \item Code-Snippets für \lstinline|socket.connect|, \lstinline|send|, \lstinline|recv|
            \item Anleitung zur Fehlerbehandlung bei Verbindungsabbrüchen
            \item Debugging-Tipps, wenn Nachrichten nicht ankommen
            \item Schritt-für-Schritt-Hilfekarten (nur bei Bedarf zugänglich)
        \end{itemize} &
        \begin{itemize}
            \item \textbf{Erste Übung:} Zwei Python-Skripte auf demselben Rechner verbinden
            \item Übung: Eine \enquote{Hallo}-Nachricht an den Server senden und Antwort empfangen
            \item Fehlerübungen: Verbindungsfehler identifizieren und beheben
            \item Mini-Tests zu IP/Port-Adressen und deren Bedeutung
        \end{itemize}
    \end{tabular}
    \caption{Didaktische Struktur der Lernaufgabe: Chat-Kommunikation}
\end{table}

\section{Learning Task 1: Lokale Verbindung via \texorpdfstring{\texttt{localhost}}{localhost}}

\begin{scaffoldBox}
    \textbf{Vorausgehende Konzepte}: Die SuS kennen bereits die historischen Hintergründe von Netzwerken, die Grundkonzepte von Netzwerkschichten, IP-Adressen, Sequenzierung und Ports, sowie die Client-Server-Architektur aufgrund einer vorangehenden Unterrichtssequenz. In der Vorausgehenden Unterrichtssequenz haben sich die SuS mit den grundlegenden Funktionen und der Bedeutung von Netzwerken auseinandergesetzt, indem Sie ein Spiel gespielt haben, bei dem sie Nachrichten über ein simuliertes Netzwerk via ausgedruckte Kärtchen zusenden mussten (s. \cref{sec:netzwerkspiel}). Dabei haben sie die Rolle von IP-Adressen und Ports kennengelernt und verstanden, wie Datenpakete in einem Netzwerk weitergeleitet werden.

    \textbf{Neue Konzepte}: Client-Server-Architektur, IP-Adresse, Port, Socket, localhost.

    Die neuen Begriffe werden zuerst anhand eines Videos (\href{https://www.youtube.com/watch?v=PBWhzz_Gn10}{Warriors of the Net}) und einer Analogie (Postsystem) eingeführt. Anschliessend wird die Grundstruktur eines Socket-Programms in Python vorgestellt. Die SuS lösen dabei einige vorbereitete Teilübungen (Part-task Practice), um die Konzepte zu verinnerlichen, und lösen danach die komplexe Lernaufgabe, bei der eine lokale Chat-Kommunikation über zwei Python-Skripte realisiert wird.
\end{scaffoldBox}

\begin{myexercise}[Chatten auf demselben Rechner]\label{ex:chat-localhost}

    \textbf{Aufgabe:} Schreiben Sie zwei Python-Programme, welche auf demselben Rechner laufen: \texttt{server\_localhost.py} und \texttt{client\_localhost.py}. Der Client verbindet sich mit dem Server und schickt eine kurze Nachricht: \enquote{Hallo Server, hier ist der Client}. Der Server empfängt die Nachricht und gibt sie auf der Konsole aus, zudem antwortet er an den Client \enquote{Hallo Client, hier ist der Server}.

    \textbf{Verwenden Sie folgende Grundstruktur:}

    \begin{lstlisting}[caption={server\_localhost.py}]
import socket

PORT = ... # HIER Port-Nummer Ihrer Wahl einsetzen

# Server-Socket erstellen
server = socket.socket(socket.AF_INET, socket.SOCK_STREAM)
server.bind(("localhost", PORT))   # lokale IP, Port 12345
server.listen(1)
print("Warte auf Verbindung...")

# Verbindung akzeptieren
conn, addr = server.accept()
print(f"Verbunden mit {addr[0]} über Port {addr[1]}")

# Nachricht empfangen
nachricht = conn.recv(1024).decode()
print(f"Empfangen: {nachricht}")

# Antwort an den Client senden
conn.send("Hallo Client, hier ist der Server!".encode())

conn.close()
server.close()
\end{lstlisting}

    \begin{lstlisting}[caption={client\_localhost.py}]
import socket

PORT = ... # HIER Port-Nummer Ihrer Wahl einsetzen

# Client-Socket erstellen und mit dem Server verbinden
client = socket.socket(socket.AF_INET, socket.SOCK_STREAM)
client.connect(("localhost", 12345))

# Nachricht an den Server senden
client.send("Hallo Server, hier ist der Client!".encode())

# Antwort vom Server empfangen
antwort = client.recv(1024).decode()
print("Antwort vom Server:", antwort)

client.close()
\end{lstlisting}
    Wählen Sie einen geeigneten Port (zwischen 1024 und 65535) und ersetzen Sie die Platzhalter in beiden Programmen.

    Lassen Sie \texttt{server\_localhost.py} auf der Konsole laufen und führen Sie dann \texttt{client\_localhost.py} aus, indem Sie in VS Code auf den Dropdown-Button rechts vom Run-Button clicken und \enquote{Run Python File in Dedicated Terminal} anklicken.

    Erweitern Sie nun Ihr Programm, indem der Client eine Nachricht schickt, und der Server antwortet mit \enquote{Nachricht erhalten!}.
    \begin{myanswer}
        \lstinputlisting[caption={server\_localhost.py}]{Grundlagen_Info/07_Netzwerke/Code/0_localhost/server_localhost.py}
        \lstinputlisting[caption={client\_localhost.py}]{Grundlagen_Info/07_Netzwerke/Code/0_localhost/client_localhost.py}
    \end{myanswer}
\end{myexercise}

\begin{myexercise}[Verständnisfragen]
    \begin{itemize}
        \item Was geschieht, wenn in \texttt{server\_localhost.py} die IP-Adresse \lstinline|"localhost"| durch \lstinline|"127.0.0.1"| ersetzt wird?
        \item Was geschieht, wenn eine andere Port-Nummer als im Client und Server verwendet wird? Funktioniert die Verbindung? Weshalb (nicht)?
        \item Warum muss der Server gestartet werden, bevor der Client versucht, sich zu verbinden?
        \item Was passiert, wenn der Server eine Nachricht sendet, aber der Client keine Antwort erwartet?
    \end{itemize}
    \begin{myanswer}
        \begin{itemize}
            \item \lstinline|"localhost"| und \lstinline|"127.0.0.1"| sind äquivalent und verweisen beide auf die lokale IP-Adresse des eigenen Rechners.
            \item Die Verbindung funktioniert nicht, da der Client versucht, sich mit einem Port zu verbinden, auf dem der Server nicht lauscht. Beide müssen denselben Port verwenden.
            \item Der Server muss zuerst gestartet werden, damit er auf Verbindungsanfragen lauschen kann. Wenn der Client zuerst startet, gibt es keinen Prozess, der die Verbindung annimmt, was zu einem Fehler führt.
            \item Wenn der Server eine Nachricht sendet, aber der Client keine Antwort erwartet, kann es sein, dass die Nachricht verloren geht oder der Client sie nicht verarbeitet, was zu Kommunikationsproblemen führen kann.
        \end{itemize}
    \end{myanswer}
\end{myexercise}

\begin{mychallenge}
    Versuchen Sie, den Client zu starten, bevor der Server läuft. Welche Fehlermeldung erhalten Sie? Wie können Sie das Problem beheben?
\end{mychallenge}


\begin{myexercise}[Reflexionsfragen (formatives Assessment)]
    Notiere nach Abschluss von Learning Task 1 (max. 5 Minuten) kurze Antworten ins Lernjournal. Tausche dich danach mit einer Person aus und ergänze fehlende Punkte.

    \textbf{Fragen:}
    \begin{enumerate}
        \item Erkläre in \emph{eigenen Worten}: Was ist ein Socket? Welche zwei Operationen hast du verwendet?
        \item Unterschied Socket vs. Port?
        \item Unterschied \gls{IP}-Adresse vs. Port (verwende die Post-Analogie).
        \item Warum muss der Server vor dem Client gestartet werden? Welche konkrete Fehlermeldung erscheint sonst?
        \item Weshalb muss \lstinline|bind()| nur einmal aufgerufen werden?
        \item Skizziere den Ablauf deiner Nachricht (Methodenaufrufe in Reihenfolge) vom Client zum Server.
        \item Nenne mindestens einen Debugging-Schritt, falls keine Antwort zurückkommt.
        \item Formuliere eine mögliche Ursache für \texttt{ConnectionRefusedError}.
        \item Kurzer Selbstcheck (Ja/Nein):
              \begin{itemize}
                  \item Ich kann den Code ohne Vorlage erneut schreiben.
                  \item Ich verstehe jede Zeile von \texttt{server.py}.
                  \item Ich weiss, wann \texttt{recv} blockiert.
              \end{itemize}
        \item Formulieren Sie ein persönliches Lernziel für Task 2 (ein Satz).
    \end{enumerate}

    \begin{myanswer}
        \begin{itemize}
            \item Socket: Endpunkt einer TCP-Verbindung; genutzt: \texttt{socket()}, \texttt{bind()}, \texttt{listen()}, \texttt{accept()}, \texttt{connect()}, \texttt{send()}, \texttt{recv()}, \texttt{close()}.
            \item Socket = Verbindungspunkt (Python-Objekt); Port = Nummer zur Identifikation eines Dienstes auf einem Rechner.
            \item \gls{IP} = Adresse des Rechners; Port = „Tür“/Dienst auf diesem Rechner (Post: Hausnummer vs. Zimmernummer).
            \item Server zuerst: Sonst kein Prozess lauscht $\rightarrow$ \texttt{ConnectionRefusedError: [Errno 61] Connection refused}.
            \item \texttt{bind()} weist Socket eine Adresse/Port zu; nur einmal nötig, da Socket danach gebunden ist.
            \item Ablauf: 
            \begin{itemize}
                \item \textbf{Client}: \texttt{socket()} $\rightarrow$ \texttt{connect()} $\rightarrow$ \texttt{send()}; 
                \item \textbf{Server}: \texttt{socket()} $\rightarrow$ \texttt{bind()} $\rightarrow$ \texttt{listen()} $\rightarrow$ \texttt{accept()} $\rightarrow$ \texttt{recv()} $\rightarrow$ (optional \texttt{send()} Antwort).
            \end{itemize}
            
            \item Debugging: Prüfen Port, laufender Server, Firewall, Print-Statements vor/nach \texttt{send/recv}, \texttt{netstat} zur Belegung.
            \item Ursache \texttt{ConnectionRefusedError}: Server nicht gestartet / falscher Port / Firewall blockiert.
            \item Blockierendes \texttt{recv}: Wartet bis Daten eintreffen oder Gegenstelle schliesst.
            \item Lernziel-Beispiel: „Ich möchte in Task 2 mehrere Nachrichten nacheinander austauschen können.“
        \end{itemize}
    \end{myanswer}
\end{myexercise}

\section{Learning Task 2: Chat-Kommunikation über ein lokales Netzwerk}

\begin{scaffoldBox}
    \textbf{Neue Konzepte}: IP-Adresse im lokalen Netzwerk, Firewall-Einstellungen.

    Die SuS lernen, wie sie die eigene IP-Adresse im lokalen Netzwerk ermitteln können (z. B. über den Befehl \texttt{ipconfig} in der Windows-Eingabeaufforderung oder \texttt{ifconfig} im Terminal auf macOS/Linux). Zudem wird erklärt, wie Firewall-Einstellungen den Datenverkehr beeinflussen können. Die Teilübungen (Part-task Practice) helfen den SuS, ihre Programme anzupassen und Verbindungsprobleme zu lösen. Am Ende steht die komplexe Lernaufgabe, bei der die Chat-Kommunikation über zwei Geräte im selben Netzwerk realisiert wird.

    \textbf{Technische Herausforderungen}: Die Lehrperson sollte vorab testen, ob die Firewall-Einstellungen auf den Schulrechnern die Kommunikation über Sockets erlauben. Gegebenenfalls müssen Ausnahmen für die verwendeten Ports eingerichtet werden oder ein eigenes Netzwerk (z. B. über einen mobilen Hotspot oder einen eigenen WLAN-fähigen Router) genutzt werden.
\end{scaffoldBox}

\begin{myexercise}[Chatten im lokalen Netzwerk]

    \textbf{Aufgabe:} Erweitern Sie Ihre Codes aus \cref{ex:chat-localhost}, so dass ihre Programme auf \textbf{verschiedenen Rechnern im selben lokalen Netzwerk} laufen. Der Server lauscht auf seiner lokalen IP-Adresse und einem Port (z.B. 12345). Der Client verbindet sich mit der IP-Adresse des Servers und schickt eine kurze Nachricht („Hallo LAN“). Der Server empfängt die Nachricht und gibt sie auf der Konsole aus.

    \textbf{Schritte:}
    \begin{enumerate}
        \item Ermitteln Sie die lokale IP-Adresse des Servers (s. \cref{ex:meine-ip}).
        \item Passen Sie die \texttt{bind}- und \texttt{connect}-Befehle entsprechend an.
        \item Starten Sie \texttt{server.py} auf Rechner A und \texttt{client.py} auf Rechner B.
    \end{enumerate}

    \textbf{Erweiterung:} Lassen Sie den Client sowie den Server mehrere Nachrichten hin- und herschicken.

    \begin{myanswer}
        \lstinputlisting[caption={\texttt{chat\_localnet\_client.py}}]{Grundlagen_Info/07_Netzwerke/Code/1_LAN/chat_localnet_client.py}
        \lstinputlisting[caption={\texttt{chat\_localnet\_server.py}}]{Grundlagen_Info/07_Netzwerke/Code/1_LAN/chat_localnet_server.py}
    \end{myanswer}
\end{myexercise}

\begin{myexercise}[Reflexionsfragen (formatives Assessment)]
    Notiere nach Abschluss von Learning Task 2 (max. 5 Minuten) kurze Antworten ins Lernjournal. Tausche dich danach mit einer Person aus und ergänze fehlende Punkte.

    \textbf{Fragen:}
    \begin{enumerate}
        \item Wie unterscheidet sich die IP-Adresse im lokalen Netzwerk von der \texttt{localhost}-Adresse?
        \item Warum ist es wichtig, die richtige IP-Adresse des Servers zu verwenden?
        \item Welche Fehlermeldung erscheint, wenn der Client versucht, sich mit einer falschen IP-Adresse zu verbinden?
        \item Nenne mindestens einen Debugging-Schritt, falls keine Verbindung hergestellt werden kann.
        \item Wie kannst du überprüfen, ob die Firewall die Verbindung blockiert?
        \item Kurzer Selbstcheck (Ja/Nein):
              \begin{itemize}
                  \item Ich verstehe jede Zeile von \texttt{server.py}.
                  \item Ich weiss, wie ich die lokale IP-Adresse meines Rechners finde.
              \end{itemize}
        \item Ein persönliches Lernziel für Task 3 (ein Satz).
        \item Welche Herausforderungen könnten auftreten, wenn Sie versuchen, über das Internet zu kommunizieren, anstatt nur im lokalen Netzwerk?
    \end{enumerate}
    \begin{myanswer}
        \begin{itemize}
            \item Unterschied: \texttt{localhost} (127.0.0.1) verweist nur auf den eigenen Rechner; die lokale IP (z.B. 192.168.x.x) identifiziert den Rechner im LAN für andere Geräte.
            \item Richtige IP nötig, sonst versucht der Client einen nicht existierenden Host oder anderen Rechner zu erreichen $\rightarrow$ keine oder falsche Verbindung.
            \item Mögliche Fehlermeldungen: \texttt{ConnectionRefusedError} (Dienst nicht aktiv), Timeout (keine Antwort), \texttt{OSError: No route to host}.
            \item Debugging: IP neu ermitteln; Server läuft? Port korrekt? \texttt{ping} testen; \texttt{netstat} / \texttt{ss} prüfen; einfache Prints vor \texttt{connect}; anderen Port probieren.
            \item Firewall prüfen: Temporär deaktivieren (nur testweise), Regel für Python-Port hinzufügen, Test mit \texttt{telnet <IP> <Port>} oder \texttt{nc}; Prüfen ob Verbindungsversuch geloggt wird.
            \item Selbstcheck: Code reproduzierbar? Jede Zeile verstanden? IP finden gelingt? Wenn Nein $\rightarrow$ gezielt nacharbeiten.
            \item Lernziel Task 3: \enquote{Ich kann einen Relay-Server nutzen, um gezielt Nachrichten an einen benannten Client zu senden.}
            \item Herausforderungen Internet: NAT und Port-Forwarding nötig; dynamische öffentliche IP; zusätzliche Latenz; Firewalls/Provider-Restriktionen; Sicherheit (Verschlüsselung, Authentifizierung); DNS-Auflösung.
        \end{itemize}
    \end{myanswer}

\end{myexercise}
\section{Learning Task 3: Chat-Kommunikation über das Internet, mittels Relay-Server}

\begin{scaffoldBox}
    \textbf{Neue Konzepte}: Öffentliche IP-Adresse, NAT, Port-Forwarding, Relay-Server, DNS (Domain Name System).

    Die SuS lernen, wie sie die öffentliche IP-Adresse ihres Netzwerks ermitteln können (z. B. über Websites wie \href{https://whatismyipaddress.com/}{whatismyipaddress.com}). Zudem wird das Konzept von NAT (Network Address Translation) und Port-Forwarding erklärt, um zu verstehen, warum direkte Verbindungen oft nicht möglich sind. Zusätzlich wird das Domain Name System (DNS) eingeführt: DNS übersetzt menschenlesbare Domain-Namen (wie \texttt{chat.example.com}) in IP-Adressen, damit Programme Server im Internet finden können. Die Teilübungen (Part-task Practice) helfen den SuS, ihre Programme anzupassen und Verbindungsprobleme zu lösen. Am Ende steht die komplexe Lernaufgabe, bei der die Chat-Kommunikation über einen Relay-Server realisiert wird.

    \textbf{Technische Herausforderungen}: Die Lehrperson sollte vorab testen, ob die Firewall-Einstellungen auf den Schulrechnern die Kommunikation über Sockets erlauben. Gegebenenfalls müssen Ausnahmen für die verwendeten Ports eingerichtet werden oder ein eigenes Netzwerk (z. B. über einen mobilen Hotspot oder einen eigenen WLAN-fähigen Router) genutzt werden.
\end{scaffoldBox}

\begin{myexercise}[Chatten über einen öffentlichen Relay-Server]

    \textbf{Aufgabe:} Sie verbinden sich mit einem bereits laufenden öffentlichen Relay-Server, der Nachrichten zwischen mehreren Clients weiterleitet. Ziel ist, dass Sie und Ihre Mitschüler:innen sich gegenseitig Nachrichten schicken können.

    \textbf{Schritte:}
    \begin{enumerate}
        \item Öffnen Sie \texttt{client.py} und passen Sie die \texttt{connect}-Zeile an:
\begin{lstlisting}[language=Python]
client.connect(("88.198.193.239", 12345))
\end{lstlisting}
        \item Senden Sie eine Nachricht (z.B. \enquote{Hallo an alle!}) an den Server.
        \item Empfangen Sie die Antwort und geben Sie sie auf der Konsole aus.
        \item Starten Sie \texttt{client.py} auf mehreren Rechnern gleichzeitig und tauschen Sie Nachrichten aus.
        \item Probieren Sie aus, ob Sie von verschiedenen Rechnern aus gleichzeitig Nachrichten senden und empfangen können. Was passiert, wenn mehrere Clients verbunden sind? Können Sie sich gegenseitig Nachrichten schicken?
        \item \textbf{Erweiterung:} Jeder Client gibt beim Start einen eigenen Namen ein. Beim Senden einer Nachricht wählen Sie aus, an welchen Namen (also an welchen anderen Client) die Nachricht geschickt werden soll. Die Nachricht wird dann nur an diesen Empfänger weitergeleitet.

    \end{enumerate}

    \textbf{Verwenden Sie folgende Grundstruktur:}

    \begin{lstlisting}[caption={client.py}]
    import socket

    client = socket.socket(socket.AF_INET, socket.SOCK_STREAM)
    client.connect(("88.198.193.239", 12345))  # IP-Adresse des Relay-Servers, Port 12345

    name = input("Dein Name: ")
    client.send(name.encode())  # Namen an Server senden

    def empfange():
        while True:
        antwort = client.recv(1024).decode()
        if not antwort:
            break
        print(antwort)

    import threading
    threading.Thread(target=empfange, daemon=True).start()

    while True:
        empfaenger = input("Empfänger-Name (oder 'exit'): ")
        if empfaenger.lower() == "exit":
        break
        nachricht = input("Nachricht: ")
        if nachricht.lower() == "exit":
        break
        # Format: EMPFAENGER:Nachricht
        client.send(f"{empfaenger}:{nachricht}".encode())

    client.close()
\end{lstlisting}
    \textbf{Hinweis:} Der Server muss die Namen der verbundenen Clients verwalten und Nachrichten entsprechend weiterleiten. So können Sie gezielt Nachrichten an einzelne Personen schicken.
    \begin{myanswer}
        \begin{lstlisting}[caption=\texttt{client.py}]
import socket
import threading

client = socket.socket(socket.AF_INET, socket.SOCK_STREAM)
client.connect(("88.198.193.239", 12345))  # IP-Adresse des Relay-Servers, Port 12345

name = input("Dein Name: ")
client.send(name.encode())  # Namen an Server senden

def empfange():
    while True:
        antwort = client.recv(1024).decode()
        if not antwort:
            break
        print(antwort)

threading.Thread(target=empfange, daemon=True).start()

while True:
    empfaenger = input("Empfänger-Name (oder 'exit'): ")
    if empfaenger.lower() == "exit":
        break
    nachricht = input("Nachricht: ")
    if nachricht.lower() == "exit":
        break
    # Format: EMPFAENGER:Nachricht
    client.send(f"{empfaenger}:{nachricht}".encode())

client.close()
        \end{lstlisting}
        \textbf{Hinweis:} Der Server muss die Namen der verbundenen Clients verwalten und Nachrichten entsprechend weiterleiten. So können Sie gezielt Nachrichten an einzelne Personen schicken.
    \end{myanswer}
    \begin{myanswer}
        \begin{lstlisting}[caption=\texttt{server.py}]
import socket
import threading

server = socket.socket(socket.AF_INET, socket.SOCK_STREAM)
server.bind(("0.0.0.0", 12345))
server.listen()

clients = []
names = []

def handle_client(client):
    while True:
        try:
            nachricht = client.recv(1024).decode()
            if not nachricht:
                break
            empfaenger, nachricht = nachricht.split(":", 1)
            if empfaenger in names:
                index = names.index(empfaenger)
                clients[index].send(f"{names[clients.index(client)]}: {nachricht}".encode())
        except:
            break

    client.close()
    clients.remove(client)
    names.remove(names[clients.index(client)])

while True:
    client, addr = server.accept()
    clients.append(client)
    name = client.recv(1024).decode()
    names.append(name)
    print(f"{name} hat sich verbunden.")
    threading.Thread(target=handle_client, args=(client,), daemon=True).start()
server.close()
            \end{lstlisting}
    \end{myanswer}
\end{myexercise}

\begin{myexercise}
    [Reflexionsfragen (formatives Assessment)]
    Notiere nach Abschluss von Learning Task 3 (max. 5 Minuten) kurze Antworten ins Lernjournal. Tausche dich danach mit einer Person aus und ergänze fehlende Punkte.
    \textbf{Fragen:}
    \begin{enumerate}
        \item Was ist der Unterschied zwischen einer lokalen IP-Adresse und einer öffentlichen IP-Adresse?
        \item Erkläre kurz, was NAT (Network Address Translation) ist und warum es in Heimnetzwerken verwendet wird.
        \item Warum ist Port-Forwarding notwendig, um Verbindungen von aussen zu einem Gerät im lokalen Netzwerk herzustellen?
        \item Was ist ein Relay-Server und welche Rolle spielt er in diesem Szenario?
        \item Wie funktioniert DNS (Domain Name System) und warum ist es wichtig für die Kommunikation im Internet?
        \item Kurzer Selbstcheck (Ja/Nein):
              \begin{itemize}
                  \item Ich kann den Code ohne Vorlage erneut schreiben.
                  \item Ich verstehe jede Zeile von \texttt{client.py}.
                  \item Ich weiss, wie ich die öffentliche IP-Adresse meines Netzwerks finde.
              \end{itemize}
        \item Reflektiere: Welche Herausforderungen könnten auftreten, wenn Sie versuchen, über das Internet zu kommunizieren, anstatt nur im lokalen Netzwerk?
    \end{enumerate}
    \begin{myanswer}
        \begin{itemize}
            \item Lokale IP: Identifiziert ein Gerät im lokalen Netzwerk (z.B.      192.168.x.x); öffentliche IP: Identifiziert das gesamte Netzwerk im Internet.
            \item NAT übersetzt lokale IP-Adressen in eine öffentliche IP-Adresse, damit mehrere Geräte im Heimnetzwerk über eine einzige öffentliche IP-Adresse kommunizieren können.
            \item Port-Forwarding leitet eingehende Verbindungen von der öffentlichen IP-Adresse an die richtige lokale IP-Adresse und den entsprechenden Port im Heimnetzwerk weiter.
            \item Ein Relay-Server empfängt Nachrichten von Clients und leitet sie an die vorgesehenen Empfänger weiter, wodurch direkte Verbindungen zwischen Clients vermieden werden.
            \item DNS übersetzt Domain-Namen in IP-Adressen, damit Menschen sich leicht zu merkende Namen verwenden können, anstatt sich numerische IP-Adressen merken zu müssen.
            \item Selbstcheck: Code reproduzierbar? Jede Zeile verstanden? Öffentliche IP finden gelingt? Wenn Nein $\rightarrow$ gezielt nacharbeiten.
            \item Herausforderungen Internet: NAT und Port-Forwarding nötig; dynamische öffentliche IP; zusätzliche Latenz; Firewalls/Provider-Restriktionen; Sicherheit (Verschlüsselung, Authentifizierung); DNS-Auflösung.
        \end{itemize}
    \end{myanswer}
\end{myexercise}